\section{Cyber-Physical Systems}\label{sec:cps}

"Cyber-Physical Systems" (CPS) are systems in which physical processes are tightly integrated with computational and communication elements. The physical parts (machines, sensors, actuators, infrastrutruce) interact with embedded computation and networking, enabling monitoring, control, and dynamic adaptation. CPS differs from more conventional embedded systems in that they are networked, often large-scale, and involve real-time feedback loops between the physical world and computation. Key early work includes Monostori et al. (2016)~\cite{monsori-cps-in-manufacturing}, which presents the foundational view of CPS / Cyber-Physical Production Systems (CPPS) as central to industry 4.0. Another core is Monostori (2014) "Cyber-physical Production Systems: Roots, Expectations and R\&D Challenges"~\cite{monsori-cpps-roots-expectations-rnd-challenges}.

CPS are seen as essential for enabling the responsiveness, flexibility, and adaptability required in \ac{hmlv}) manufacturing, because they allow systems to sense context, respond in close to real time, and reconfigure or adjust behavior when parts, batches, or processes change.

\subsection{Role of CPS in Manufacturing}

In manufacturing, CPS integrates production equipment with computation and connectivity, enabling continuous monitoring optimization, and reconfiguration. For \ac{hmlv} contexts, where product variation and small batch sizes are common, CPS make it possible to:

\begin{itemize}
    \item Collect and analyze sensor data for real-time decision making.
    \item Enable predictive maintenance and reduce unplanned downtime.
    \item Support dynamic scheduling and resource allocation under fluctuating demand.
    \item Improve quality assurance through online defect detection.
\end{itemize}

Lee et al. (2015)~\cite{LEE-2015-a-cps-for-indsutry40-based-manufacturing-systems} proposed the well-known 5C architecture (connection, conversion, cyber, cognition, configuration), which structures how data from shop-floor devices are transformed into actionable insights. Such architectures illustrate the layered integration of the physical and cyber domains, which ultimately enables self-adaptive production systems.

\subsection{Framework and Tools}

Several frameworks have been proposed to standardize CPS in manufacturing. The Reference Architectural Model Industry 4.0 (RAMI 4.0)~\cite{Schweichhart2015RAMI4.0} provides a three dimensional model linking hierarchy levels, product lifecycle stages, and architecture layers, situating CPS within the broader Industry 4.0 ecosystem.

Research at Aarhus University, led by Larsen et al.~\cite{larsen-et-al-into-cps}, has contributed to the development of INTO-CPS, a co-simulation platform for model-based CPS engineering. INTO-CPS supports the design and validation of CPS by integrating multiple modeling tools and enabling virtual experiments across domains. Such platforms exemplify how academic research provides concrete tools for industry to adpot CPs in practice.

\subsection{Surveys and Applications}

Several reviews summarize the state of CPS in manufacturing:

\begin{itemize}
    \item Dafflon et al. (2021)~\cite{dafflon-the-challenges-approaches} reviewed 78 CPS studies (2010-2019), identifying common challenges such as interoperability, scalability, and human integration, and highlighting technologies such as Internet of Things (IoT), Big Data, and Artificial Intelligence (AI).
    \item Liu et al. (2018)~\cite{liu-cps-based-smart-warehouse} focused on CPS-based smart warehouses, showing how CPS concepts extend to logistics, with applications in sensor-driven monitoring, automated transport, and adaptive storage systems.
    \item Benzadri et al. (2021)~\cite{benzadri-a-modelling-framework-for-cps-based-industry} proposed a modeling framework for CPS-based manufacturing systems, demonstrating the practical challenges of applying CPS in industrial settings.
\end{itemize}

These studies show that CPS research is both broad and maturing, with applications ranging from production lines to logistics and supply chain contexts.

\subsection{Challenges and Future Directions}

Despite progress, several open issues remain:

\begin{itemize}
    \item Interoperability between heterogeneous machines and software platforms.
    \item Scalability of CPS solutions for Small and Medium-sized Enterprises (SMEs) and \ac{hmlv} contexts.
    \item Data quality, security, and privacy in highly connected environments.
    \item Human-in-the-loop integration, ensuring that CPS augment rather than replace operators (Industry 5.0).
\end{itemize}

Looking ahead, CPS are increasingly converging with digital twin technologies, enabling tighter synchronization between physical assets and virtual models. In industry 5.0, CPS are expected to evolve toward human-centric and sustainable architectures, where adaptability is balanced with transparency, explainability, and resource efficiency.

